% =====================================================================
%  Annexe D : Donnees de demonstration
% =====================================================================
\chapter{Données de démonstration}
\label{ann:donnees}

\chapitreresume{Pour tester le catalogue sans tout saisir à la main, cette annexe
fournit une commande de gestion qui crée des auteurs, des catégories, des livres,
des membres et des emprunts. Elle crée aussi deux comptes de démonstration et
explique comment réinitialiser les données.}

\section{La commande de démonstration}

Créez l'arborescence :

\begin{codebox}[tree]{Arborescence de la commande}
catalogue/
`-- management/
    |-- __init__.py
    `-- commands/
        |-- __init__.py
        `-- donnees_demo.py
\end{codebox}

\begin{codebox}[python]{catalogue/management/commands/donnees\_demo.py}
from datetime import date, timedelta

from django.contrib.auth.models import Group, User
from django.core.management.base import BaseCommand
from django.utils import timezone

from catalogue.models import (
    Auteur,
    Categorie,
    Emprunt,
    Livre,
    Membre,
)

AUTEURS = [
    ("Antoine", "Saint-Exupéry", 1900),
    ("Mariama", "Bâ", 1929),
    ("Albert", "Camus", 1913),
    ("Chinua", "Achebe", 1930),
    ("Ahmadou", "Kourouma", 1927),
]

CATEGORIES = [
    ("Roman", "roman", "Fictions longues et récits."),
    ("Jeunesse", "jeunesse", "Albums et romans pour les jeunes."),
    ("Poésie", "poesie", "Recueils et textes poétiques."),
    ("Afrique", "afrique", "Littératures africaines."),
    ("Histoire", "histoire", "Essais et ouvrages historiques."),
]

LIVRES = [
    (
        "Le Petit Prince",
        "le-petit-prince",
        "9782070612758",
        "Antoine",
        "Saint-Exupéry",
        1943,
        3,
        ["Jeunesse", "Roman"],
        "Un aviateur rencontre un petit prince venu d'une autre planète.",
    ),
    (
        "Une si longue lettre",
        "une-si-longue-lettre",
        "9782264030305",
        "Mariama",
        "Bâ",
        1979,
        2,
        ["Roman", "Afrique"],
        "Une lettre qui raconte la vie de femmes sénégalaises.",
    ),
    (
        "L'Étranger",
        "l-etranger",
        "9782070360024",
        "Albert",
        "Camus",
        1942,
        2,
        ["Roman"],
        "Meursault raconte un fait divers sous le soleil d'Alger.",
    ),
    (
        "Le Monde s'effondre",
        "le-monde-s-effondre",
        "9782264030831",
        "Chinua",
        "Achebe",
        1958,
        1,
        ["Roman", "Afrique"],
        "Le destin d'Okonkwo dans un village igbo en mutation.",
    ),
    (
        "Les Soleils des indépendances",
        "les-soleils-des-independances",
        "9782020028010",
        "Ahmadou",
        "Kourouma",
        1968,
        1,
        ["Roman", "Afrique"],
        "Fama, prince déchu, traverse les indépendances.",
    ),
]


class Command(BaseCommand):
    help = "Crée un jeu de données de démonstration pour GesiBiblio."

    def add_arguments(self, parser):
        parser.add_argument(
            "--reset",
            action="store_true",
            help="Supprime les emprunts et les livres de démonstration.",
        )

    def handle(self, *args, **options):
        if options["reset"]:
            self.reset()

        auteurs = {}
        for prenom, nom, naissance in AUTEURS:
            auteur, _ = Auteur.objects.get_or_create(
                prenom=prenom,
                nom=nom,
                defaults={
                    "date_naissance": date(naissance, 1, 1),
                },
            )
            auteurs[nom] = auteur

        categories = {}
        for nom, slug, description in CATEGORIES:
            categorie, _ = Categorie.objects.get_or_create(
                slug=slug,
                defaults={"nom": nom, "description": description},
            )
            categories[nom] = categorie

        for (
            titre,
            slug,
            isbn,
            prenom,
            nom,
            annee,
            exemplaires,
            noms_categories,
            resume,
        ) in LIVRES:
            livre, cree = Livre.objects.get_or_create(
                slug=slug,
                defaults={
                    "titre": titre,
                    "isbn": isbn,
                    "auteur": auteurs[nom],
                    "annee_publication": annee,
                    "nombre_exemplaires": exemplaires,
                    "resume": resume,
                },
            )
            if cree:
                livre.categories.set(
                    [categories[n] for n in noms_categories]
                )
                self.stdout.write(f"  + {livre.titre}")

        membres = self.creer_membres()
        self.creer_emprunts(membres)

        self.stdout.write(
            self.style.SUCCESS("Données de démonstration prêtes.")
        )

    def creer_membres(self):
        membres = []
        for numero, username, prenom, nom in [
            ("M00001", "lecteur1", "Awa", "Ndiaye"),
            ("M00002", "lecteur2", "Kofi", "Mensah"),
        ]:
            user, cree = User.objects.get_or_create(
                username=username,
                defaults={
                    "first_name": prenom,
                    "last_name": nom,
                    "email": f"{username}@example.com",
                },
            )
            if cree:
                user.set_password("Demonstration2026!")
                user.save()

            membre, _ = Membre.objects.get_or_create(
                user=user,
                defaults={"numero": numero},
            )
            membres.append(membre)

        lecteurs, _ = Group.objects.get_or_create(name="Lecteurs")
        for membre in membres:
            membre.user.groups.add(lecteurs)

        return membres

    def creer_emprunts(self, membres):
        aujourdhui = timezone.localdate()
        livre = Livre.objects.filter(slug="le-petit-prince").first()
        if not livre or not membres:
            return

        Emprunt.objects.get_or_create(
            livre=livre,
            membre=membres[0],
            date_emprunt=aujourdhui - timedelta(days=5),
            defaults={
                "date_retour_prevue": aujourdhui + timedelta(days=9),
            },
        )

        livre2 = Livre.objects.filter(slug="l-etranger").first()
        if livre2:
            Emprunt.objects.get_or_create(
                livre=livre2,
                membre=membres[1],
                date_emprunt=aujourdhui - timedelta(days=30),
                defaults={
                    "date_retour_prevue": aujourdhui - timedelta(days=2),
                },
            )

    def reset(self):
        Emprunt.objects.all().delete()
        Livre.objects.filter(
            slug__in=[livre[1] for livre in LIVRES]
        ).delete()
        self.stdout.write("Livres et emprunts de démonstration supprimés.")
\end{codebox}

\begin{terminal}
python manage.py donnees_demo
\end{terminal}

\begin{terminal}
  + Le Petit Prince
  + Une si longue lettre
  + L'Étranger
  + Le Monde s'effondre
  + Les Soleils des indépendances
Données de démonstration prêtes.
\end{terminal}

\begin{note}
La commande est idempotente : la relancer ne crée pas de doublons grâce à
\code{get\_or\_create}. Le mot de passe des comptes de démonstration est volontaire
et documenté. Ne l'utilisez jamais en production.
\end{note}

\section{Les comptes créés}

\begin{table}[htbp]
\centering
\small
\begin{tabularx}{\textwidth}{L{3.4cm} L{4.4cm} X}
\toprule
\textbf{Utilisateur} & \textbf{Mot de passe} & \textbf{Rôle} \\
\midrule
\code{lecteur1} & \code{Demonstration2026!} & Membre Awa Ndiaye, groupe
Lecteurs. \\
\code{lecteur2} & \code{Demonstration2026!} & Membre Kofi Mensah, groupe
Lecteurs. \\
\code{admin} & défini à la création & Superutilisateur, tous les droits. \\
\bottomrule
\end{tabularx}
\caption{Comptes de démonstration.}
\end{table}

\begin{attention}
Changez immédiatement ces mots de passe si vous déployez le projet sur un serveur
accessible. Un mot de passe de démonstration est public par construction.
\end{attention}

\section{Vérifier les données}

\begin{codebox}[python]{Shell Django}
from catalogue.models import Auteur, Categorie, Emprunt, Livre, Membre

Auteur.objects.count()        # 5
Categorie.objects.count()     # 5
Livre.objects.count()         # 5
Membre.objects.count()        # 2
Emprunt.objects.count()       # 2

Livre.objects.filter(categories__slug="afrique")
Emprunt.objects.filter(date_retour__isnull=True)
\end{codebox}

\section{Réinitialiser}

\begin{terminal}
python manage.py donnees_demo --reset
python manage.py donnees_demo
\end{terminal}

\begin{attention}
\code{--reset} supprime les emprunts et les livres dont le slug figure dans la
liste de démonstration. Il ne supprime ni les auteurs, ni les catégories, ni les
membres. Pour repartir de zéro, utilisez une base de développement dédiée :
\code{python manage.py flush} vide toutes les données, mais conserve le schéma.
\end{attention}

\section{Alternative : une fixture JSON}

Pour un jeu de données figé, Django accepte des fixtures. Créez le fichier
\code{categories.json} dans un dossier \code{fixtures/} de l'application :

\begin{codebox}[json]{catalogue/fixtures/categories.json}
[
  {
    "model": "catalogue.categorie",
    "pk": 1,
    "fields": {
      "nom": "Roman",
      "slug": "roman",
      "description": "Fictions longues et récits."
    }
  },
  {
    "model": "catalogue.categorie",
    "pk": 2,
    "fields": {
      "nom": "Jeunesse",
      "slug": "jeunesse",
      "description": "Albums et romans pour les jeunes."
    }
  }
]
\end{codebox}

\begin{terminal}
python manage.py loaddata categories
\end{terminal}

\begin{note}
La commande \code{loaddata} cherche le fichier dans les dossiers
\code{fixtures/} des applications installées. Une fixture est pratique pour des
données de référence stables, moins pour des données qui doivent être recalculées
ou créées avec des mots de passe hachés.
\end{note}

\begin{bonne}
Pour les données de référence (catégories, langues, pays), une migration de données
ou une commande idempotente est préférable à une fixture. Pour des données de test
figées, la fixture reste simple et lisible.
\end{bonne}
